\begin{corollary}[Positive Euler number and cohomology nonvanishing]
\hypertarget{tag-03IV}{}\label{wbII:test:cdp-comparison}
\StatementRecord{wbII:test:cdp-comparison}{proved}{03IV}{\SeriesRef{wbII:thm:corrected-cdp-global}, \SeriesRef{wbII:cor:corrected-cdp-relative}, \SeriesRef{wbII:cor:candidate-conductor-actual}; \citet[Section~10.3, Theorem~10.5 and Corollary~10.6, pp.~88--89]{AC2026}}

For the standard compact threefold $X$, every $L\in\operatorname{Pic}(X)$
satisfies
\[
 H^2(X,T_X\otimes L)\ne0,
 \qquad R^2f_*(T_X\otimes L)\ne0.
\]
The second conclusion also follows directly from the cusp conductor.
\end{corollary}


\begin{proof}
The standard compact smooth threefold constructed in Paper I has
$H^1(X;\mathbb Z)=H^2(X;\mathbb Z)=0$, holomorphic algebraic reduction
to $\mathbb P^1$, and $\langle c_3(T_X),[X]\rangle=2$.
If the displayed degree-two cohomology vanished for one line bundle,
\cref{wbII:thm:corrected-cdp-global} would give $c_3(X)\le0$, a contradiction.
If the second direct image vanished for one twist,
\cref{wbII:cor:corrected-cdp-relative} would give such a global twist after
tensoring by the pullback of a sufficiently positive base line bundle,
again a contradiction. Independently,
\cref{wbII:cor:candidate-conductor-actual} constructs a nonzero cusp-fibre
cohomology class and identifies it with the fibre of the second direct
image by top-degree base change. The resulting counterexample to the CDP conclusion is due to \citet[Section~10.3, Theorem~10.5 and Corollary~10.6, pp.~88--89]{AC2026}.

\end{proof}